\subsection{Historical Supervision Experiments: Evidence and Limits}
\label{app:supervision-audit}

Completed training, negative performance, inconclusive evidence and unavailable supervision are different outcomes. Retained training logs establish that a DistilRoBERTa mention--slot classifier was actually fitted, rather than merely implemented. Its weak labels selected the first numerically matching mention. The largest retained run (job 854626) used 230/50/50 internal train/development/test instances drawn from the existing NLP4LP evaluation corpus. The logged selection counts are 51/280 slots and 0/45 legacy exact-instance outcomes. Five nominal test instances do not enter that latter denominator. Those metrics are not the all-required-scalar criterion used in this manuscript.

A later summary's claim that this run was worse than a rule baseline on every metric conflicts with the preserved stdout. The audit also identifies an evaluator path that can instantiate an untrained neural head for a nominal rule configuration. Missing original run checkpoints, exact run source SHAs and prediction traces prevent resolving that comparison by replay. Therefore these logs demonstrate completed weak-label learning with poor recorded absolute outcomes, but not a clean causal comparison of neural learning against rules. No headline result is substituted from the conflicting summary.

P0 is a separate completed classifier-plus-assignment study. Logistic regression over 25 handcrafted features and one frozen MiniLM cosine feature was selected on development slot selection; the encoder was not fine-tuned. On its 50 oracle-schema cases, independent selection has coarse readiness 0.80, exact maximum-weight assignment 0.78, and assignment plus repair 0.78, versus a separately retained typed-greedy comparator at 0.86. These are the saved readiness outcomes, not whole-instance numerical correctness. The probability outputs have not been shown to be calibrated; the exact assignment has no private null option, and several components differ from the comparator. Its first-value-match labels and evaluation-origin split do not establish an independent semantic test.

Both histories used internal subsets of the 331-query corpus; disjoint example IDs alone did not establish independence. The saved text-group audit finds train/development and development/test overlap. This compromises any later use of those fits as independent evidence on the full 331 queries. It does not retroactively turn every internal observation into a nonexistent experiment. The defensible lesson is narrower: fitting an equality-derived target and optimizing its scores did not establish trustworthy occurrence-specific supervision.

MiniLM cross-encoder and multitask grounder implementations also exist, but a completed non-smoke run for the former and a completed NL4Opt auxiliary neural checkpoint for the latter were not verified. They are not counted as negative model results. Actual GAMS and targeted-synthetic auxiliary runs used only two and one effective test problems; they are too small and weakly supervised to justify general transfer claims. Those execution records remain in the audit inventory rather than a new performance table. Fine-tuned retrieval addresses a different task, and an encoder substitution alone is not a new relationship-grounding contribution.

\subsection{NL4Opt Relationship and Independence Audit}
\label{app:nl4opt-audit}

The audit reads 713 NL4Opt source problems at pinned revision \texttt{49f1e0d66b7fdcd33305a7f281c2a7c13f5620ea}. Its 3,114 positive proposals and 3,246 ambiguous links are derived occurrence--declaration candidates. They do not constitute new independent human annotations. A conservative source-template investigation propagates quarantine to 399 problems; 314 remain unflagged but uncertified. Earlier lexical-screen negatives and later structural-template concerns are compatible because they test different relations. Quarantine counts are precautionary grouping outcomes, not measured contamination or semantic-error rates.

A separate relationship-labeling recovery audit is discussed in the main paper. Its numerical provider comparison, recovery protocol and complete-case denominators appear in main Section~\ref{M-subsec:relationship-audit} and Online Resource 1, Section~\ref{app:relationship-evidence}. No model consensus has been accepted as a human semantic reference, and no model was trained from that packet in the reported analyses.
